% Source: https://artofproblemsolving.com/community/c3h1335926_function_equation
% Stable identity: aops-topic-1335926-q01
Putting $a=b=0$ gives $2f(0)=f(0)$, so $a_0=f(0)=0$.

For fixed $b$ with $|b|<10$, the polynomial $g(a)=f(a+b)-f(a)-f(b)$ vanishes for all $a$ in the interval $(-10,10)$, which has infinitely many points, so $g$ is the zero polynomial. The same then holds for every $b$, so $f(a+b)=f(a)+f(b)$ is a polynomial identity in $a$ and $b$. Its terms of total degree $k\ge2$ are $a_k\bigl((a+b)^k-a^k-b^k\bigr)$, and $(a+b)^k-a^k-b^k$ contains the term $ka^{k-1}b\ne0$. Hence $a_k=0$ for all $k\ge2$, and $f(x)=a_1x$.

Then $a_n+\cdots+a_0=f(1)=a_1=324$, and all other coefficients are $0$. The list $a_{n-1},a_{n-3},\dots,a_3,a_1$ is therefore $0,0,\dots,0,324$ (it contains $a_5=a_3=0$ because $n\ge6$). As a set, $S=\{0,324\}$, whose average is
\[
\dfrac{0+324}{2}=\boxed{162}.
\]
(If repeated zeros were counted separately, the average $\dfrac{648}{n}$ would depend on $n$; reading $S$ as a set makes the answer independent of $n$.)
